\section{Regularization and exact stationary matching}\label{regularization}
The noise formula of Definition~2.46 and Proposition~2.47
in~\cite{odonnell} makes smoothing a natural candidate operation.
Smoothing the potential, however, does not commute with its nonlinear
stationary-density construction. The following statements separate a
preserved inequality from a failed exact equality.

\begin{proposition}\label{smoothingfloor}
Let $h>0$ satisfy $\E hPh=1$ and $hPh\geq b>0$. For $0\leq\tau\leq1$ put
\[
 H_\tau=P_\tau h,\qquad Z_\tau=\E H_\tau P H_\tau,\qquad
 h_\tau=H_\tau/\sqrt{Z_\tau},\qquad \pi_\tau=h_\tau P h_\tau.
\]
Then $\pi_\tau\geq b$ for every $\tau$. No marginal-matching or moment
assumption is needed for this assertion.
\end{proposition}
\begin{proof}
The noise kernels commute. Cauchy--Schwarz in each probability row of
$P_\tau$ therefore gives
\[
 H_\tau P H_\tau=(P_\tau h)P_\tau(Ph)
 \geq\bigl(P_\tau\sqrt{hPh}\bigr)^2\geq b.
\]
By Walsh Parseval,
\[
 0<Z_\tau=\sum_S\rho^{|S|}\tau^{2|S|}\widehat h(S)^2
 \leq\sum_S\rho^{|S|}\widehat h(S)^2=1.
\]
The constant-character multiplier is one, including at $\tau=0$.
Dividing the pointwise inequality by $Z_\tau$ proves the result.
\end{proof}

\begin{theorem}\label{smoothingmatchingfailure}
Fix $0<\rho<1$, $0<b<1$, $L\geq3$ and $1\leq k\leq L-2$, and put
$s=(1+\rho)/(1-\rho)$. There is a positive potential $h$ with
$\E hPh=1$ such that its stationary density has exact uniform labelled
marginals through order $k$, while
\[
 \pi_\tau\geq b,\qquad \E\pi_\tau^s\leq\sqrt2<2
       \quad(0\leq\tau\leq1).
\]
For one specified coordinate $j$, the same potential satisfies
$\E\pi_\tau\chi_{\{j\}}\ne0$ for every $0<\tau<1$.
Thus every nontrivial smoothing loses even a first marginal while the
stationary floor and critical cap survive for this constructed family.
\end{theorem}
\begin{proof}
Choose $A\subseteq[L]$ with $m=|A|=k+1$ and $j\notin A$, and put
$B=A\cup\{j\}$, $\alpha=\rho^m$ and $\beta=\rho^{m+1}$.
Fix the dimension-independent number
\[
 \epsilon=\min\left\{\frac18,\frac{1-b}{4},
                         \frac{\log2}{64s}\right\}>0
\]
and define $w=1+\epsilon(\chi_A+\chi_B)$.
This density has mean one and
$w\geq1-2\epsilon\geq\max\{3/4,(1+b)/2\}>b$.
Its only nonconstant Walsh coefficients are at $A$ and $B$, whose
degrees exceed $k$, so it has the required labelled marginals.
We next realize it as the stationary density of an actual positive
potential rather than treating the density and potential as independent.

Under uniform input, $y=\chi_A$ and $z=x_j$ are independent uniform
signs. Since their coordinate sets are disjoint, Boolean $\rho$-noise
induces on $(y,z)$ the exact product kernel
$Q=P_\alpha\otimes P_\rho$. In particular
\[
 P[f(\chi_A,x_j)]=(Qf)(\chi_A,x_j)
\]
for every function $f$ on these four points. The eigenvalues of $Q$ are
$1,\alpha,\rho,\beta$; all entries are strictly positive.
Apply Lemma~\ref{inversestationary} on this four-point space to the path
$w_t=1+t\epsilon(y+yz)$, $0\leq t\leq1$. It is bounded below by $1/2$.
The unique positive solutions $g_tQg_t=w_t$ start from $g_0=1$ and obey
\[
 \|\log g_1\|_\infty
 \leq\frac{\lceil\tfrac12\log_2(4/(1/2))\rceil+2}{2(1/2)}
                  \|\epsilon(y+yz)\|_\infty
 =8\epsilon=:\eta\leq\frac{\log2}{8s}.
\]
The inverse estimate uses four points, not $2^L$, and is uniform as
$\alpha$ becomes small. The lift $h=g_1(\chi_A,x_j)$ consequently
satisfies $hPh=w$, $\E hPh=1$, and $e^{-\eta}\leq h\leq e^\eta$.

Every $H_\tau=P_\tau h$ and $PH_\tau$ lies in the same interval, so
$e^{-2\eta}\leq Z_\tau\leq e^{2\eta}$ and
$\pi_\tau\leq e^{4\eta}$. Thus
$\E\pi_\tau^s\leq e^{4s\eta}\leq\sqrt2$ for all $\tau$.
Proposition~\ref{smoothingfloor} supplies the floor for all $\tau$.
At $\tau=1$ these are the original stationary density and its cap.

Write the unique four-character expansion
\[
 h=c+d\chi_{\{j\}}+e\chi_A+f\chi_B,\qquad c=\E h>0.
\]
Equating the three indicated coefficients in $hPh=w$ gives
\begin{align}
 (1+\rho)cd+(\alpha+\beta)ef&=0,\label{smoothingcancel}\\
 (1+\alpha)ce+(\rho+\beta)df&=\epsilon,\label{smoothingcoefA}\\
 (1+\beta)cf+(\rho+\alpha)de&=\epsilon.\label{smoothingcoefB}
\end{align}
If $e=0$, equation~\eqref{smoothingcancel} forces $d=0$, contradicting
\eqref{smoothingcoefA}. If $f=0$, the same equation again forces $d=0$,
contradicting \eqref{smoothingcoefB}. Hence $ef\ne0$.
Now
\[
 H_\tau=c+\tau d\chi_{\{j\}}+\tau^m e\chi_A
                         +\tau^{m+1}f\chi_B.
\]
Using \eqref{smoothingcancel} in the $\chi_{\{j\}}$ coefficient of
$H_\tau PH_\tau$ proves the exact identity
\begin{equation}\label{smoothingdefect}
 \E\pi_\tau\chi_{\{j\}}
   =\frac{(\alpha+\beta)ef}{Z_\tau}
                         (\tau^{2m+1}-\tau).
\end{equation}
For $0<\tau<1$ this is nonzero, since $m\geq2$, $ef\ne0$ and
$\alpha+\beta,Z_\tau>0$. No perturbative remainder or numerical
nonzero test is used. At $\tau=0$ the normalized density is uniform,
and at $\tau=1$ it is the original exactly matched density.
\end{proof}

For fixed $C>0$, the required $k_L=\min\{L,\lfloor C\sqrt L\rfloor\}$
eventually belongs to $[1,L-2]$. Apply the theorem in those dimensions
and use $h=1$ in the finitely many remaining dimensions. This gives the
claimed obstruction within the original growing matching regime.
No uniform lower bound on the magnitude of \eqref{smoothingdefect} is
asserted: $\alpha$ tends to zero as $k_L$ grows. Exact matching is the
equality that fails. The constructed potentials have a fixed positive
minimum, so they are not a collapsing-gap family.

The distinction between potential and density regularization is essential.
For $w_\tau=P_\tau\pi$, stochasticity preserves mean one and the floor,
Jensen's inequality and invariance of $\mu$ give
$\E w_\tau^s\leq\E\pi^s$, and the Walsh formula preserves every zero
low-degree coefficient. Lemma~\ref{inversestationary} then supplies the
unique actual positive potential $u_\tau$ with $u_\tau Pu_\tau=w_\tau$.
These elementary invariances do not compare the gaps of $K_h$ and
$K_{u_\tau}$. A dimension-uniform spectral comparison under the original
hypotheses remains unproved; it cannot be replaced by assuming a positive
minimum for the reconstructed potential or a lower bound on $Ph$.
